% Термины и определения. Порядок — по алфавиту (кириллица, затем латиница).
% Нумерация id произвольная, но уникальная.

\newglossaryentry{t-encoding}{
    name={кодирование колонки},
    description={преобразование последовательности однотипных значений столбца в более компактное двоичное представление с использованием статистических свойств данных (повторов, ограниченного диапазона, общих префиксов и т. п.)}
}
\newglossaryentry{t-skipping}{
    name={пропуск данных},
    description={(англ. data skipping) приём оптимизации чтения, при котором по статистикам (минимум, максимум, число null-значений) блоков данных заранее определяется, что блок не содержит подходящих под предикат строк, и его чтение пропускается}
}
\newglossaryentry{t-compaction}{
    name={слияние (компакция)},
    description={(англ. compaction) фоновая операция в LSM-дереве, объединяющая несколько отсортированных файлов одного уровня в меньшее число файлов следующего уровня с удалением перезаписанных и удалённых записей}
}
\newglossaryentry{t-format}{
    name={файловый формат},
    description={спецификация двоичного представления табличных данных в файле, включающая разбиение на блоки, схему хранения значений (построчное или поколоночное), способы кодирования, сжатия и встроенные метаданные (статистики, индексы)}
}
\newglossaryentry{t-tableformat}{
    name={открытый табличный формат},
    description={(англ. open table format) спецификация метаданных поверх набора файлов данных, обеспечивающая представление их как единой таблицы с поддержкой транзакций, эволюции схемы, моментальных снимков и параллельного доступа; примеры — Apache Iceberg, Delta Lake, Apache Hudi, Apache Paimon}
}
\newglossaryentry{t-columnar}{
    name={колоночный формат хранения},
    description={формат, в котором значения одного столбца таблицы располагаются в файле непрерывно; обеспечивает высокую степень сжатия однотипных данных и чтение только запрашиваемых столбцов}
}
\newglossaryentry{t-predicate}{
    name={предикат},
    description={логическое условие фильтрации строк (например, \texttt{col > 100 AND col < 200}); может передаваться в файловый формат для отсечения блоков данных на уровне декодера (predicate pushdown)}
}
\newglossaryentry{t-mor}{
    name={слияние при чтении},
    description={(англ. merge-on-read) стратегия таблиц с первичным ключом, при которой обновления и удаления записываются как новые версии, а актуальное состояние строки вычисляется во время чтения объединением версий из всех уровней LSM-дерева по полю последовательности}
}
\newglossaryentry{t-lakehouse}{
    name={lakehouse},
    description={архитектура хранения данных, объединяющая дешёвое масштабируемое хранилище объектов озера данных (data lake) с транзакционными гарантиями, управлением схемой и SQL-доступом, характерными для хранилищ данных (data warehouse)}
}
\newglossaryentry{t-lsm}{
    name={LSM-дерево},
    description={(англ. Log-Structured Merge-tree) структура данных, оптимизированная под интенсивную запись: новые данные буферизуются в памяти и сбрасываются на диск отсортированными неизменяемыми файлами (уровень~0), которые затем фоновым слиянием перемещаются на нижележащие уровни}
}
\newglossaryentry{t-olap}{
    name={OLAP},
    description={(англ. Online Analytical Processing) класс нагрузок, характеризующихся сложными аналитическими запросами с агрегацией над большими объёмами данных и редкими изменениями; противопоставляется OLTP}
}
\newglossaryentry{t-oltp}{
    name={OLTP},
    description={(англ. Online Transaction Processing) класс нагрузок, характеризующихся большим числом коротких транзакций чтения и записи отдельных записей по ключу}
}
